 \documentclass[11pt]{article}
\usepackage{amsmath,amsfonts}
\usepackage{graphicx}
%\usepackage{xcolor}
%\definecolor{gray}{rgb}{.75,.75,.75}
%\usepackage[landscape]{geometry}
%\usepackage{hyperref}
%\usepackage[bookmarks=true, pdfborder={0 0 .5}, urlbordercolor=gray]{hyperref}

\renewcommand{\thefootnote}{\fnsymbol{footnote}}


\addtolength{\topmargin}{-.75in}
\addtolength{\textheight}{1.5in}
\addtolength{\oddsidemargin}{-.35in}
\addtolength{\textwidth}{.7in}
\pagestyle{empty}

\newcommand{\ds}{\displaystyle}
\newcommand{\ts}{\textstyle}

\newcommand{\Z}{\mathbb{Z}}
\newcommand{\R}{\mathbb{R}}
\newcommand{\C}{\mathbb{C}}
\newcommand{\EE}{\mathcal{E}}
\newcommand{\tZ}{\tilde{\Z}}

\newcommand{\Sym}{\mathrm{Sym}}

\renewcommand{\circle}[1]{{\bigcirc\hspace{-.75em}#1}}

\begin{document}

\footnote[0]{Notes by Peter Magyar \ {\tt magyar@math.msu.edu}}
\vspace{-1em}

\centerline{\bf Math 133\hspace{1.5em} \hfill{\Large\bf
Parametric Calculus}\hfill 
Stewart \S10.2}
\vspace{1em}

\noindent 
{\bf Tangents of a parametric curve.} We have learned how to write a curve parametrically, as the path of a particle whose position at time $t$ is
given by two coordinate functions $(x(t),y(t))$ over a time interval $t\in[a,b]$.

Considering the curve as a track on which the particle runs,
the {\it tangent line} at a point $(x(c),y(c))$ 
is the path the particle would take if
it were suddenly released from the track at time $t=c$, 
keeping a constant velocity from that moment.
The velocity at $t=c$ has horizontal and vertical components
$(x'(c), y'(c))$, giving the parametric line:
$$
(x(c)+x'(c)\,t, \, y(c) +y'(c)\,t)\,.
$$
This is the line which best approximates the curve near the point of tangency $(x(c),y(c))$, with the time set so that the line passes through the point at $t=0$.

We can convert this parametric line into an $xy$-equation as in \S10.1.
The slope is the horizontal over the vertical velocity: $m=\frac{y'(c)}{x'(c)}$, and we know the line passes through $(x(c),y(c))$, so we have 
the point-slope equation:
$$
y = \tfrac{y'(c)}{x'(c)}(x{-}x(c))+ y(c)\,.
$$
Here $(x,y)$ is a general point of the line, 
but $x(c),y(c), x'(c),y'(c)$ are constants 
computed from the coordinate functions of the original curve.

To further explain this, 
we imagine the original curve as the graph of a function
$y=f(x)$, meaning $y(t)=f(x(t))$ for all $t$. 
The Chain Rule gives:
$$
y'(t)  \ =\ f'(x(t))\cdot x'(t) 
\qquad\Longleftrightarrow\qquad
\frac{dy}{dt} \ =\ \frac{dy}{dx}\cdot\frac{dx}{dt}\ .
$$
At time $t=c$ and $x=x(c)$, this gives our previous slope formula:
$$
f'(x(c))=\frac{y'(c)}{x'(c)}
\qquad\Longleftrightarrow\qquad
\frac{dy}{dx}  \ =\ 
\frac {\ \ \frac{dy}{dt}\ \ } {\frac{dx}{dt}}\,.
 $$ 
{\bf Tangents of a circle.} We find the tangent line to 
$(x(t),y(t)) = (2\sin(\pi t),2\cos(\pi t))$ at the point $(\sqrt 2,\sqrt 2)$.
First, to picture the curve, we note:
\begin{itemize}
\item Since the components are $2\sin$ and $2\cos$ of the 
same quantity, the curve is a circle of radius 2.
\item The full circle is traced by $\pi t\in[0,2\pi]$, i.e.~$t\in[0,2]$.
\item The curve starts at $(x(0), y(0)) = (0,2)$
on the $y$-axis; it moves clockwise, since the $x$-coordinate
$2\sin(\pi t)$ increases for small $t\geq0$.

\end{itemize}

\centerline{
\includegraphics[width=2.8in]
{10-2-Circle.png} 
}
\vspace{1em}

\noindent
To apply our formulas, we need to know the value $t=c$ at which
the curve passes through the given point:
$(x(t),y(t))=(\sqrt 2,\sqrt 2)$. That is, we must solve the 
system of equations:
$$\left\{
\begin{array}{c}
2\sin(\pi t) = \sqrt 2\\
2\cos(\pi t) = \sqrt 2
\end{array}
\right.
\qquad\Longleftrightarrow\qquad
t = \tfrac14\,.
$$
We can find a simultaneous solution to both equations precisely because the point
lies on the curve. We have $(x'(c),y'(c)) = (2\pi\cos(\frac\pi4), -2\pi\sin(\frac\pi4)) = (\sqrt 2\pi, -\sqrt 2\pi)$,
so the tangent line is:
$$
(x(c)+ x'(c)t, y(c) + y'(c)t) \ =\ (\sqrt 2 + \sqrt 2\pi t, \sqrt 2 -\sqrt2\pi t)
$$
$$
y \ =\ 
\tfrac{y'(c)}{x'(c)}(x{-}x(c))+ y(c)
\ =\ 
\tfrac{\sqrt 2\pi}{-\sqrt 2\pi}(x{-}\sqrt 2) + \sqrt 2
\qquad\Longleftrightarrow\qquad
y = -x+2\sqrt2\,.
$$
Note that each tangent to the circle is perpendicular to the corresponding radius.
\\[1em]
{\bf Tangents of a polynomial curve.} Find the tangent to $(x(t),y(t))=(t^2, t^3-3t)$ at the point $(3,0)$.
This is not a familiar curve, so to picture it, we must plot points by plugging in various
values of $t$:
\\[1em]
\centerline{
\includegraphics[width=3in]
{10-2-Poly.png} 
}


\noindent
We see that the curve passes twice through the given point $(3,0)$. Algebraically:
$$
\left\{
\begin{array}{c}
t^2 = 3\\
t^3-3t = 0
\end{array}
\right.
\qquad\Longleftrightarrow\qquad
t = \sqrt3 \ \ \text{or}\ \ t = -\sqrt3\,.
$$
Note that $t^3-3t=0$ by itself has the solutions $t=0,\pm\sqrt3$, but
$t=0$ does not satisfy the first equation $t^2=3$:
for time $t=0$, the curve is at $(0,0)$, not $(3,0)$.

Now we can easily find the two tangent lines: $(3+2\sqrt 3 t,  6t)$ and 
$(3-2\sqrt 3 t,  6t)$.
\\[.5em]
{\sc example:} Which points of this curve have horizontal tangents?
The tangent is horizontal when the vertical velocity is zero: $(t^3-3t)' = 3t^2-3 =0
\ \Longleftrightarrow\ t=\pm1$, corresponding to the points $(1, -2)$ and $(1,2)$.
\\[1em]
{\bf Arclength.}  After applying derivatives to parametric curves, we now apply integrals,
which compute the size or bulk of geometric objects. The most natural measure of the size of a curve is its arclength. We already computed this for graph curves $y=f(x)$ in \S8.1, 
and now we do the more general parametric case.  

%We recall our general scheme for computing any measure of size $S$, which might be
%length, area, volume, work, etc.
%We cut the geometric object into slices whose position is determined by
%some variable $x\in[a,b]$. 
%We mark off the interval $[a,b]$ into $n$ increments of width $\Delta x$,
%each with a sample point $x_i$\,.
%This splits the object into $n$ slices, 
%and summing up their sizes gives the total size:  $S=\sum_{i=1}^n S_i$.
%Because the slice at $x_i$ is so thin, 
%we can find a good approximation of its size by a simple formula
%$\Delta S_i \approx f(x_i)\Delta x$.
%Taking $n\to\infty$ and $\Delta x\to 0$, the approximations become exact, 
%and we have:
%$$
%S \ =\ \lim_{n\to\infty} \sum_{i=1}^n f(x_i)\Delta x \ =\ \int_a^b\!\! f(x)\,dx\,.
%$$

We follow the general scheme for computing any measure of size of a geometric object
from \S5.2.
We want the arclength $L$ of a parametric curve $(x(t),y(t))$ for $t\in[a,b]$. 
We cut the curve into $n$ bits determined by $\Delta t$-increments of $t\in[a,b]$.
\\[1em]
\centerline{
\includegraphics[width=3.6in]
{10-2-Length.png} 
}
\vspace{-.5em}

\noindent
Because the bit at the sample point $t_i$ is so short, 
it is well approximated by a straight segment, and we can 
use the Pythagorean Theorem to compute its length:
$$
\Delta L_i  \ \approx\ \sqrt{(\Delta x)^2+(\Delta y)^2}
\ =\ 
\sqrt{\frac{(\Delta x)^2+(\Delta y)^2}{(\Delta t)^2}}\  \Delta t
\ =\ 
\sqrt{(\tfrac{\Delta x}{\Delta t})^2+(\tfrac{\Delta y}{\Delta t})^2}\  \Delta t\,.
$$
In the limit as $n\to \infty$, we get $\Delta t\to0$  and $\frac{\Delta x}{\Delta t}\to \frac{dx}{dt} = x'(t_i)$; similarly for $\frac{\Delta y}{\Delta t}$:
$$
L \ =\ \lim_{n\to\infty} \sum_{i=1}^n\Delta L_i \ =\ \lim_{n\to\infty}\sum_{i=1}^n\! \sqrt{(\tfrac{\Delta x}{\Delta t})^2+(\tfrac{\Delta y}{\Delta t})^2}\ \Delta t
\ = \ \int_a^b\!\!\! \sqrt{(\tfrac{dx}{dt})^2+(\tfrac{dy}{dt})^2}\ dt\,.
$$
In Newton notation:
$$
L  \ = \ \int_a^b\!\!\! \sqrt{x'(t)^2+y'(t)^2}\ dt\,.
$$
In fact, the integrand is just the total speed of the particle at time $t$, combining the horizontal and vertical speeds. To be precise, $L$ is the {\it total distance traveled} along the curve;
if the trajectory repeats its motion around a closed loop, or if it changes direction and backtracks, 
this is larger than the geometric length of the curve.
\\[.5em]
{\sc example:} Compute the circumference length of a circle of radius $r$.
The standard parametrization is $(x(t),y(t))=(r \cos(t), r\sin(t))$ for $t\in[0,2\pi]$, with 
derivative $(x'(t),y'(t)) = (-r\sin(t),r\cos(t))$, and length:
$$
L \ =\ \int_0^{2\pi} \!\!\!\!\sqrt{(-r\sin(t))^2 + (r\cos(t))^2}\,dt
\ =\ \int_0^{2\pi} \!\!r \sqrt{\sin^2(t) + \cos^2(t)}\,dt
$$
$$
\ =\ \int_0^{2\pi} \!\!r \,dt \ =\ \left.\vphantom{\sum}rt\,\right|_{t=0}^{t=2\pi} \ =\ 2\pi r\,.
$$
The integral is so easy because the particle travels at constant speed $r$.
This was much harder in \S8.1, using our previous formula $L = \int_{-r}^r\! \sqrt{1{+}f'(x)^2}\ dx$, where $f(x)=\sqrt{r^2-x^2}$.
\\[.5em]
{\sc example:} Find the length of one arch of the cycloid from \S10.1:
$(x(t),y(t))=(t{-}\sin(t), 1{-}\cos(t))$ for $t\in[0,2\pi]$. We have $(x'(t),y'(t))=$ $(1{-}\cos(t),\, \sin(t))$, so:
$$
L
\ =\ 
\int_0^{2\pi}\!\!\! \sqrt{(1{-}\cos(t))^2+(\sin(t))^2}\,dt
\ =\ 
\int_0^{2\pi}\!\!\! \sqrt{1{-}2\cos(t)+\cos^2(t)+\sin^2(t)}\,dt
$$
$$
\ =\  
\int_0^{2\pi}\!\!\! \sqrt{2(1{-}\cos(t))}\,dt 
\ =\
\int_0^{2\pi}\! 2\sin(\tfrac t2)\,dt\ =\ 8\,.
$$
Here we used the identity $\sin(\frac t2) =\sqrt{\frac{1-\cos(t)}2}$\,.
%\\[1em]
%{\bf Area under a parametric curve.}
%$A =\int_{x(a)}^{x(b)} y\,dx 
%=\int_a^b y(t)\,\frac{dx}{dt}\,dt$
%\\
%area of circle $\pi r^2$, cycloid $3\pi = 3(\text{area of wheel})$.
\bigskip

\noindent
{\bf Area.}  Consider a parametric curve $(x(t),y(t))$ for $t\in [a,b]$
which closes into a loop starting and ending at the point $(x(a),y(a))=(x(b),y(b))$.
Then the enclosed area is:
$$
A \ = \ \pm \int_a^b y(t) x'(t)\,dt \ =\ \mp \int_a^b x(t) y'(t)\,dt,
$$
where the signs after the two equalities are $+,-$ if the loop travels clockwise around the region, 
but $-,+$ if it travels  counterclockwise.

As an exercise, prove these formulas by slice analysis. 
The point is that area is added while traveling in one direction,
then subtracted while traveling back.


\end{document}

%%%%%%%%%%%  Picture  %%%%%%%%%%
\\[1em]
\centerline{
%\includegraphics[width=2in]
{1-8-Discont-iv.png}
\qquad \qquad 
%\includegraphics[width=2in]
{1-8-Discont-v.png} 
}
\vspace{0em}

%%%%%%%%%%%%%%%%%%%%%%%%%

%%%%%%%%%%%  Table  %%%%%%%%%%
In fact, we get the following derivatives:
$$\begin{array}{|c||c|c|c|c|c|c|}
\hline &&&&&& \\[-.9em] 
f(x) & \sin(x) & \cos(x) & \tan(x)  & \sec(x) & \csc(x) & \cot(x)
\\[.2em] \hline &&&&&& \\[-.9em]
f'(x) &\cos(x) & -\sin(x) & \sec^2(x) & \tan(x)\sec(x) & -\cot(x)\csc(x) & -\csc^2(x)
\\[.1em] \hline
\end{array}$$
\\[1em]
%%%%%%%%%%%%%%%%%%%%%%%%%%



















